\begingroup
\renewcommand{\arraystretch}{1.45}
\InterpolationArticle{19}
Particularly frequent, and therefore deserving the greatest attention, is the case in which
the values of the arc \(x\) for which the values of the function \(X\) are given form an arithmetic progression whose common difference equals the quotient obtained
by dividing the whole circumference into as many parts as there are terms. Thus, if
the progression were continued beyond its last term, its first terms, increased by a whole circumference, would appear again in succession. Having completed the principal points of the general investigation above, we will now treat this case, in which
the given terms constitute, as it were, a complete \emph{period}, separately and more fully, to crown the discussion. Before we begin this investigation itself, however, two
lemmas must be given first.

\textsc{First Lemma.} \emph{If the arcs \(a,b,c,d\) etc., whose number is \(\mu\), form
an arithmetic progression in which the difference of the terms \(b-a=c-b=d-c\)
etc. is \(=\frac{360^{0}}{\mu}\), the product of \(\mu\) factors \(P=\)}
\SourceLeaf{303}
\[
\begin{aligned}
&\sin\frac12(t-a)\sin\frac12(t-b)\sin\frac12(t-c)\sin\frac12(t-d)\ \mathrm{etc.}
\end{aligned}
\]
\noindent\emph{becomes \(=\mp\frac{\sin\frac12\mu(t-a)}{2^{\mu-1}}\), where the upper or lower sign applies according as \(\mu\) is even or odd.}

\emph{Proof.} This theorem can indeed readily be derived from another given by the distinguished \textsc{Euler} in \emph{Introd. in Anal. Infin.} I, \S. 240; so that nothing may seem to be missing here, however, we will briefly explain another proof.

\Needspace{30\baselineskip}
I. If the product \(P\) is first put in the form
\[
\begin{aligned}
&\cos\frac12(t-a-180^{0})\cos\frac12(t-b-180^{0})\\[5pt]&\qquad\cos\frac12(t-c-180^{0})\cos\frac12(t-d-180^{0})\mathinner{.\,.\,.}
\end{aligned}
\]
\noindent and then expanded into a sum of cosines by the known rules for sines, it will readily be seen that a function of either the form \(F\) or the form \(G\)
(Article 13) results, extending to the \(\frac12\mu^{\mathrm{th}}\) order, and that its \emph{highest} term becomes
\[
\begin{aligned}
&=\frac{1}{2^{\mu-1}}\cos\frac12(\mu t-a-b-c-d-\mathrm{etc.}-\mu\times180^{0})\\[5pt]
&=\frac{1}{2^{\mu-1}}\cos\bigl(\frac12\mu(t-a-360^{0})+90^{0}\bigr)\\[5pt]
&=-\frac{1}{2^{\mu-1}}\sin\frac12\mu(t-a-360^{0})\\[5pt]
&=-\frac{1}{2^{\mu-1}}\sin\frac12\mu(t-a)\cos\mu\times180^{0}\\[5pt]
&=\mp\frac{1}{2^{\mu-1}}\sin\frac12\mu(t-a)
\end{aligned}
\]
\noindent with the upper sign applying to an even value of \(\mu\), the lower to an odd value.

\Needspace{11\baselineskip}
II. Since \(\sin\frac12\mu(t-a)\) becomes \(=0\) when \(t\) is assigned one of the values \(a.\ b,c,d\) etc.,
the theorem of Article 14 shows that \(\sin\frac12\mu(t-a)\) is divisible by the product \(P\);
the quotient, however, must necessarily be a fixed quantity \(=h\), since \(\sin\frac12\mu(t-a)\) and \(P\)
are of the same order. But it follows evidently (by I) from the equation \(P=\frac1h\sin\frac12\mu(t-a)\)
that \(h=\mp2^{\mu-1}\), and therefore \(P=\mp\frac{\sin\frac12\mu(t-a)}{2^{\mu-1}}\)\quad Q.~E.~D.

\Needspace{32\baselineskip}
\textsc{Second Lemma.} \emph{If the arcs \(a,b,c,d\mathinner{.\,.\,.}\) have the same relation as in the first lemma, the product of \(\mu\) factors}
\[
\begin{aligned}
&(\cos t-\cos a)(\cos t-\cos b)(\cos t-\cos c)(\cos t-\cos d)\mathinner{.\,.\,.}
\end{aligned}
\]
\noindent\emph{becomes \(=\frac{\cos\mu t-\cos\mu a}{2^{\mu-1}}\)}

\emph{Proof.} Since
\SourceLeaf{304}
\[
\begin{aligned}
&\cos t-\cos a=2\sin\frac12(t-a)\sin\frac12(-t-a)\\[5pt]
&\cos t-\cos b=2\sin\frac12(t-b)\sin\frac12(-t-b)\\[5pt]
&\cos t-\cos c=2\sin\frac12(t-c)\sin\frac12(-t-c)\\[5pt]
&\cos t-\cos d=2\sin\frac12(t-d)\sin\frac12(-t-d)\\[5pt]
&\mathrm{etc.}
\end{aligned}
\]
\noindent the product of these factors becomes, by the first lemma,
\[
\begin{aligned}
&=2^{\mu}\times\mp\frac{\sin\frac12\mu(t-a)}{2^{\mu-1}}\times\mp\frac{\sin\frac12\mu(-t-a)}{2^{\mu-1}}\\[5pt]
&=\frac{\sin\frac12\mu(t-a)\sin\frac12\mu(-t-a)}{2^{\mu-2}}=\frac{\cos\mu t-\cos\mu a}{2^{\mu-1}}\quad\mathrm{Q.\ E.\ D.}
\end{aligned}
\]

\Needspace{19\baselineskip}
Since
\[
\begin{aligned}
&\sin\frac12\mu(t-a)=-\sin\frac12\mu(t-b)\\[5pt]&=\sin\frac12\mu(t-c)=-\sin\frac12\mu(t-d)\ \mathrm{etc.}
\end{aligned}
\]
\noindent and also \(\cos\mu a=\cos\mu b=\cos\mu c=\cos\mu d\) etc., it is evident that the product
in the first lemma also becomes
\[
\begin{aligned}
&=\pm\frac{\sin\frac12\mu(t-b)}{2^{\mu-1}}\\[5pt]&=\mp\frac{\sin\frac12\mu(t-c)}{2^{\mu-1}}=\pm\frac{\sin\frac12\mu(t-d)}{2^{\mu-1}}\ \mathrm{etc.}
\end{aligned}
\]
\noindent and that the product in the second lemma also becomes
\[
\begin{aligned}
&=\frac{\cos\mu t-\cos\mu b}{2^{\mu-1}}=\frac{\cos\mu t-\cos\mu c}{2^{\mu-1}}\\[5pt]&=\frac{\cos\mu t-\cos\mu d}{2^{\mu-1}}\ \mathrm{etc.}
\end{aligned}
\]

\endgroup
